\paragraph{Model merging: constructions without limits.}
A wave of post-hoc merging methods followed Task
Arithmetic~\citep{ilharco2023task}: Fisher-weighted
averaging~\citep{matena2022fisher}, sparsification approaches
(TIES~\citep{yadav2023ties}, DARE~\citep{yu2024dare},
DELLA~\citep{della2024}), rotation- and subspace-alignment families
(KnOTS~\citep{stoica2025knots}, TSPA~\citep{tspa2025},
DO-Merging~\citep{domerging2025}, Core~Space~\citep{panariello2025core},
TARA~\citep{tara2025}, ARM~\citep{arm2026streaming}), and merge-time
quantisation (TVQ~\citep{kim2025tvq},
1-bit-Merging~\citep{bit1merging2025}). All are \emph{constructions}:
each proposes an encoder and evaluates it empirically. What none of them
supplies, and what this paper does supply, is a way to tell in advance
which of them will work on a given cohort. The authors do also prove a
rate-distortion lower bound at a storage budget, but report it as
provenance rather than as a contribution, and
\S\ref{subsec:floor-zero} and \S\ref{subsec:w3} say exactly how little
of it survives contact with real adapters.

\paragraph{Existing theoretical analyses of merging.}
\emph{Weight disentanglement}~\citep{ortizjimenez2023disentanglement}
identifies the rank-$r$ Hessian structure of LoRA fine-tunes and
provides the $H_t = U_t D_t U_t^\top$ formalism the authors adopt, but gives no
rate-distortion bounds. Task-Vector Bases~\citep{jang2025taskvectorbases}
posits a thin-Gaussian-shell geometry; ATM~\citep{atm2024alternating}
and its successor~\citep{gradients2025taskvectorsgradients} develop a
second-order Taylor bound on task-arithmetic error. None of these frames
merging as lossy source coding.

\paragraph{Concurrent work using the same lens.} \citet{cao2026collapse},
posted while this work was in progress and independently of it, study
``merging collapse'' in fine-tuned LLMs, identify representational
incompatibility between tasks as a strong correlate of it, and use
rate-distortion theory to argue for a dimension-dependent limit on
mergeability that holds regardless of the merging method. The authors therefore
make no priority claim for the lens itself.

\textbf{The two bounds constrain different objects, and the authors correct one
of their own descriptions of the difference.} An earlier draft distinguished
the bound of this paper as bounding worst-task rather than average distortion. That is not a difference:
their Theorem~1 bounds $\delta_{\max}(\hat\theta) = \max_i
\E_X\|h(X;\hat\theta) - h(X;\theta_i)\|_2^2$, which is also a maximum over
tasks. The real differences are these. Theirs measures distortion in the
\emph{hidden states} of a fixed layer and the authors' in \emph{parameter} space under
a per-task curvature $H_t$, so their quantity needs forward passes and data
while the authors' is read off the adapter files. Theirs assumes linear mode
connectivity, that every convex combination of the fine-tuned minima attains
the same loss, and derives $\delta_{\max} \geq \tfrac14\Delta^2$ with a convex
merge achieving $\tfrac{d}{2(d+1)}\Delta^2$ by Jung's theorem, where $\Delta$
is the diameter of the task representation clusters. The authors' assumes rank-$r$
structure and, for the rate term only, Assumption~\ref{asm:epr}. And the role
of the rate differs sharply: their rate-distortion statement is that
$R(D) = 0$ exactly when $D \geq \Delta^2/4$ and $R(D) \geq \log_2 N$ below it,
a zero-rate threshold rather than a tradeoff, whereas the content of the authors' is
the $R$-dependence itself.

\textbf{Their central negative result deserves a direct answer, because it
looks like a counterexample to the authors'.} They report that parameter-space
conflict metrics correlate minimally with merging collapse while
representational incompatibility correlates strongly, and this paper's
headline instrument is a parameter-space measurement. The authors think the two agree
where they overlap, for a reason the authors' own null supplies. Their bound is proved
for convex merges under LMC, and convex merges are precisely the class in
which the authors \emph{also} find parameter-space geometry to have no detectable
effect: fifteen of twenty base-by-method cells in NLL
(\S\ref{subsec:ridge}) and zero of forty in downstream accuracy. The authors' positive
claim is confined to merges that solve a linear system, which are not convex
combinations of the fine-tuned models and which their assumption excludes.
Read that way, both papers find that parameter-space quantities do not predict
what averaging-type merges do; the authors add that one parameter-space quantity does
predict what solve-based merges do. The authors have not tested their $\Delta$
on the cohorts measured here, and doing so would need the forward passes this
paper's audit avoids, so the reconciliation is offered as an argument and not
as a measurement. A reader
interested in the limits of merging should read both.

\paragraph{Closed-form merging lineage.} The authors' encoder sits in a line of
closed-form merges that invert a per-task curvature surrogate.
RegMean~\citep{jin2023regmean} solves a least-squares system in the
per-task input-activation Grams (data-dependent, forward passes
required); LoRM~\citep{lorm2024} carries the closed form over to LoRA
modules in federated continual learning, and
RegMean++~\citep{regmeanpp2025} adds cross-layer dependencies to
RegMean's objective so that merged earlier layers are accounted for
when merging deeper ones. The authors' differs in two ways: it is
\emph{data-free}, since the projector comes from the adapter factors
alone, and it is derived from a worst-task rate-distortion lower bound
rather than from an empirical fit. What it shares with that line is more
important here than what separates it. All of these methods solve a linear
system in a curvature surrogate, and all carry a ridge term whose role
\S\ref{subsec:ridge} is about. The authors test RegMean directly on both cohorts
(\S\ref{subsec:regmean}) and find the same conditioning effect there, which
is what licenses talking about the family rather than about the encoder
introduced here.
LoRM and RegMean++ are not run: RegMean++'s distinguishing mechanism is a
cross-layer dependency in the activation Gram and LoRM's setting is
federated continual learning, so neither has a data-free adapter-only form
the authors could implement without inventing the part that makes it that method.
The family claim therefore rests on two members of it, not five.

\paragraph{Rate-distortion and rotation-for-quantisation ancestry.}
Classical rate-distortion theory~\citep{shannon1959rd,coverthomas}
covers the single-source setting; multi-description
coding~\citep{elgamalcover1982md} treats \emph{average}, not maximum,
distortion, which Lemma~\ref{lem:id} shows is strictly weaker. The
direct ancestor of the authors' achievability construction is
TurboQuant~\citep{zandieh2025turboquant}, whose Theorem~3
Theorem~\ref{thm:achv} recovers at $T = 1$. Structured rotations as
quantisation pre-conditioners (QuIP~\citep{tseng2024quip},
QuaRot~\citep{ashkboos2024quarot}, SpinQuant~\citep{liu2025spinquant},
RoLoRA~\citep{xi2024rolora}) motivate the authors' randomised orthogonal mixers
but are single-source.

\paragraph{A recent negative result, and why a lower bound still
matters.} A large-scale empirical study~\citep{systematic2025merging}
reports that across four LLM families only Task Arithmetic reliably
produces constructive interference, with subspace-alignment methods
often failing. Their result is an \emph{upper bound} on specific
algorithms; the authors' Theorem~\ref{thm:general} is a \emph{lower bound} on
every algorithm. The authors' measurements bear on theirs, but less directly
than an earlier reading suggested. The authors find $\deff = Tr$ in every
layer of all four bases and in both initialisation conditions, so the
task subspaces are indeed linearly independent and the irreducible floor
is zero throughout (\S\ref{subsec:floor-zero}). The authors previously took that
to mean subspace-alignment methods have no overlap to exploit, and cited
their own KnOTS measurement as support. The authors withdraw both. Formal
independence is not the absence of overlap: under shared initialisation
the subspaces are near-collinear and $\bar H$ is ill-conditioned
(\S\ref{sec:regimes}), which is exactly the situation an alignment
method might be expected to address. And the authors' KnOTS implementation
reduced algebraically to Task Arithmetic under its default inner
combination, so it never tested the question (\S\ref{sec:discussion}).
The authors make no claim here about published KnOTS. Where the authors' ordering differs
from theirs, the
difference is one of metric and setup (average interference across
$16$ heterogeneous benchmarks vs.\ worst-task NLL excess on a
controlled matrix), not a contradiction.
